\section{Discussion and conclusions}
\label{sec: discussion}

In this work we have utilized a set of magneto-hydrodynamical simulations of MW-LMC analogues from the Auriga project~\cite{Grand:2016mgo} to study the effect of the LMC on the local DM  distribution and explore its implications for DM direct detection. We first identified  15 MW-LMC candidate systems by requiring that the stellar mass of the LMC analogue and its distance from the host at its first pericenter approach match observations. We then focused on one MW-LMC analogue and studied the impact of the LMC on the local DM distribution at different times (snapshots) in its orbit. In particular, we considered four representative snapshots: the isolated MW analogue, the first pericenter approach of the LMC analogue, the closest snapshot to the present day MW-LMC system, and the MW-LMC system at a future point in time, $\sim 175$~Myr after the present day. 



We extracted the DM density and velocity distribution in the Solar region. 
The allowed positions of the Sun in the simulations were chosen such that they match the observed Sun-LMC geometry. In particular, we first found the stellar disk orientations in the simulations that make the same angle with the orbital plane of the LMC analogues as in observations. We then determined the position of the Sun for each allowed disk orientation by matching the angles between the orbital angular momentum of the LMC and the Sun's position and velocity vectors in the simulations to their observed values. The \emph{best fit Sun's position} was then defined as the one that leads to the closest match of the angles between the Sun's velocity vector and the LMC's position and velocities with observations. Using the local DM velocity distributions extracted from the simulations, we computed the halo integrals and showed how the LMC impacts their high speed tails. Finally, we simulated the signals in three near future xenon, germanium, and silicon direct detection experiments, considering the DM-nucleus interaction in the first two experiments and the DM-electron interactions in the latter, and studied the implications of the LMC on their exclusion limits. We summarize our findings below:

\begin{itemize}

\item The percentage of the DM particles originating from the LMC  in the Solar region is in the range of $[0.0077-2.8]$\% for the selected MW-LMC analogues. The local DM density is in the range of $[0.21-0.60]$~GeV/cm$^3$, depending on the halo.

\item The local speed distribution of the DM particles originating from the LMC peaks at the high speed tail ($\gtrsim 500$~km/s with respect to the  center of the MW analogue) of the local speed distribution of the native DM particles of the MW, with large halo-to-halo variations in the results. Focusing on different snapshots of one halo shows that the LMC impacts the high speed tail of the local DM speed distribution not only at its pericenter approach and the present day, but also up to $\sim 175$~Myr after the present day.

\item The LMC causes a shift in the high speed tail of the halo integrals towards larger speeds. Three key factors contribute to the variations in the tails of the halo integrals, quantified with the metric $\Delta \eta$ (eq.~\eqref{eq: delta eta}). First, a higher percentage of the DM particles originating from the LMC in the Solar region in general leads to a higher $\Delta \eta$, across different MW-LMC analogues and different snapshots of one system. Second, the exact Sun-LMC geometry for the choice of the Sun's position in the simulations 
has an impact on $\Delta \eta$, with the best fit Sun's position being close to a position which maximizes $\Delta \eta$. Therefore, in the MW we expect $\Delta \eta$ to be close to its maximum value at the Solar circle. Third, 
the native DM particles of the MW in the Solar region are boosted in response to the passage of the LMC, causing a further increase in $\Delta \eta$. The combination of  this boost and the presence of the high speed LMC particles in the Solar region causes a shift of greater than $\sim 150$~km/s in the high speed tail of the halo integrals at the present day.

\item The differences in the high speed tail of the halo integrals due to the LMC lead to considerable shifts in the expected direct detection exclusion limits  towards lower cross sections and smaller DM masses.  
In particular, the LMC lowers the exclusion limits set by the future xenon experiment on the DM-nucleon cross section by an order of magnitude for a DM mass of $\sim 8$~GeV, by more than three orders of magnitude for a DM mass of $\sim 6$~GeV, and by more than five orders of magnitude for a DM mass of $\sim 5$~GeV. For the future germanium experiment, the exclusion limits are lowered by an order of magnitude for a DM mass of $\sim 0.5$~GeV and by more than three orders of magnitude for a DM mass of $\sim 0.4$~GeV. The LMC also lowers the exclusion limits set by the future silicon experiment on the DM-electron cross section by up to a factor of $\sim 4$ for a DM mass of $\sim 1$~MeV and by up to a factor of $\sim 50$ for a DM mass of $\sim 0.6$~MeV. Furthermore, the LMC leads to a horizontal shift in the exclusion limits towards smaller DM masses, by a few GeV for xenon, a few hundred MeV for germanium, and a fraction of MeV for silicon, with the shift being more prominent for smaller DM masses. Thus, the LMC extends the parameter space probed by direct detection experiments towards lower DM masses.


\end{itemize}


The novel finding of our work is that the LMC's influence on the local DM distribution is significant even in a fully cosmological simulation, which follows the evolution of the MW and LMC analogues. While there are important halo-to-halo variations in  the results of our cosmological simulations,  a number of key conclusions could be reached by focusing on different snapshots of a particular MW-LMC analogue. Our study shows that a massive satellite that is just past its pericentric approach can significantly boost the high speed tail of the local DM velocity distribution. We also find that our particular Sun-LMC geometry maximizes the impact on the DM velocity distribution.

Our results are in general agreement with those of ref.~\cite{Besla:2019xbx}, which studied the effect of the LMC on direct detection signals in a suite  of idealized simulations of the LMC's orbit around the MW. Similar to our findings, they found that for small DM masses the LMC causes a vertical shift of more than an order of magnitude in the exclusion limits on the DM-nucleon cross section towards smaller cross sections. 


The results of our fully cosmological simulations provide further evidence of the importance of the LMC's impact on the local DM distribution. It also strengthens the argument that these significant effects should not be overlooked in the analysis of future DM direct detection data, especially for low DM masses. Finally, our results have wider implications for the validity of utilizing the idealized simulations to understand other phenomena, such as the predictions for a DM wake induced by the LMC in the halo. Future cosmological simulations, which can achieve higher resolution would ultimately be able to quantify with high precision the differences in the high speed tail of the local DM velocity distribution due to the presence of the LMC.


%*******************************%

\acknowledgments

We thank Gianfranco Bertone for discussions on the results of this work. AS, NR, and NB acknowledge the support of the Natural Sciences and Engineering Research Council of Canada (NSERC), funding reference number RGPIN-2020-07138, and the NSERC Discovery Launch Supplement, DGECR-2020-00231. AF is supported by a UKRI Future Leaders Fellowship (grant no MR/T042362/1). GB acknowledges support from NSF CAREER award AST-1941096. CSF acknowledges support from the European Research Council (ERC)
Advanced Investigator grant DMIDAS (GA 786910). FAG acknowledges support from ANID FONDECYT Regular 1211370 and by the ANID BASAL project FB210003. FAG acknowledges funding from the Max Planck Society through a ``Partner Group'' grant. RG acknowledges financial support from the Spanish Ministry of Science and Innovation (MICINN) through the Spanish State Research Agency, under the Severo Ochoa Program 2020-2023 (CEX2019-000920-S). This work used the DiRAC Memory Intensive system at Durham University, operated by ICC on behalf of the STFC DiRAC HPC Facility (www.dirac.ac.uk). This equipment was funded by BIS National E-infrastructure capital grant ST/K00042X/1, STFC capital grant ST/H008519/1, and STFC DiRAC Operations grant ST/K003267/1 and Durham University. DiRAC is part of the National E-Infrastructure.
